\documentclass[10pt,a4paper]{amsart}
\usepackage[margin=19mm]{geometry}
\usepackage{amsmath,amssymb,mathtools}
\usepackage[scr=rsfs,cal=euler]{mathalfa}
\usepackage{xfrac}
\usepackage[unicode,hidelinks,bookmarks=false]{hyperref}
\newcommand{\ie}{i.e.\ }
\newcommand{\cf}{cf.\ }
\newcommand{\resp}{respectively\ }
\newcommand{\Subsection}{\subsection}
\providecommand{\nobreakdash}{\nobreak\hbox{-}}
\setlength{\emergencystretch}{2em}
\multlinegap0pt
\usepackage{enumerate}
\usepackage{amssymb}
\usepackage[normalem]{ulem}
\newtheorem{theorem}{Theorem}
\newtheorem{remark}{Remark}
\title{Co-Higgs bundles of Schwarzenberger type and the determinant morphism}
\author{KUNTAL BANERJEE}
\date{Source: arXiv:2509.03773v2 (CC BY 4.0)}
\begin{document}
\maketitle

\noindent\emph{Source and licence.} Co-Higgs bundles of Schwarzenberger type and the determinant morphism. Version arXiv:2509.03773v2 (CC BY 4.0), distributed by arXiv under the Creative Commons Attribution 4.0 International licence, http://creativecommons.org/licenses/by/4.0/, as recorded in the arXiv licence metadata for this exact version and shown on the abstract page https://arxiv.org/abs/2509.03773v2. Full licence notice: \texttt{licenses/arxiv-2509.03773-CC-BY-4.0.txt}.\par\smallskip

\noindent\textit{Editorial scope.} Counted unit: the article's theorem det2 (source0.tex lines 306--312). The article's complete original proof of it is delivered with it (source0.tex lines 314--367, one \texttt{proof} environment of 54 lines). No mathematics is written, repaired or re-derived: the statement and the proof are the article's own lines, in the article's own notation, and no retained line is substituted or re-flowed.  No further source-local context is needed: every cross-reference of every retained line resolves inside the two delivered spans.  Nothing was omitted from any retained span, so the package carries no omission record.  No retained line contains a citation, so the wrapper carries no bibliography and no bibliography entry.  No retained line contains a figure environment, an image-inclusion command or drawing code, so no image is delivered and the package declares no asset dependency.  Editorial formatting only, in addition: an AMS \texttt{amsart} preamble that restates the article's own declarations and macros used by the retained lines, namely the environments \texttt{theorem}, \texttt{remark} on separate counters, together with the standard packages those lines need.  The retained environments print in this document as Theorem 1, Remark 1 in order of appearance, whereas the article prints the same items with the numbering of its own full document.  The complete native source of the article as distributed in the arXiv e-print for this exact version is bundled unchanged as \texttt{sources/184439/source0.tex}.
\begin{theorem}\label{det2}
There is a bijection 
\begin{align*}
&\emph{Im}\left(\det|_{\mathcal{M}_{\mathbb{P}^2}(V_1^\rho, T_{\mathbb{P}^2})}\right)\to\\
&\left(\frac{H^0(\mathcal{O}(2))^\times\times H^0(T_{\mathbb{P}^2}(-1))^\times}{(q, C)\sim (\alpha^2\cdot q, \alpha^{-1}\cdot C)}\sqcup\frac{H^0(T_{\mathbb{P}^2})}{\{\pm 1\}}\right)\bigg/[(q, C)]\sim [A]\iff q\otimes\emph{Sym}^2(C) = \emph{Sym}^2(A)
\end{align*} given by $q\otimes\emph{Sym}^2(C)\mapsto [(q, C)]$ and $\emph{Sym}^2(A)\mapsto [A]$.
\end{theorem}

\begin{proof}
Since $V_1^\rho = \mathcal{O}\oplus\mathcal{O}$ is Gieseker semistable every traceless co-Higgs bundle $\phi$ modelled on $V_1^\rho$ is Gieseker semistable. Write 
\begin{equation}
\phi = \begin{bmatrix}
  A & B\\
    C & -A
\end{bmatrix},~~ A, B, C\in H^0(T_{\mathbb{P}^2}). 
\end{equation}
Let $B, C\neq 0$. The integrability condition of $\phi$ implies that either 
\begin{equation}
\phi = \begin{bmatrix}
    \lambda & \mu\\
    \mu' & -\lambda
\end{bmatrix}\otimes C',~~ \lambda, \mu, \mu'\in H^0(\mathcal{O}(1)), C\in H^0(T_{\mathbb{P}^1}(-1))
\end{equation}
or 
\begin{equation}
\phi = \begin{bmatrix}
    \lambda & \mu\\
    1 & -\lambda
\end{bmatrix}\otimes C ,~~\mu\neq 0, \lambda\in\mathbb{C}.
\end{equation}
In the second case, $\det(\phi) = -(\lambda^2 + \mu)\mbox{Sym}^2(C)$. \\

On the other hand let $B = 0$ or $C = 0$. Without loss of generality $C=0$. So $\phi = \begin{bmatrix}
    A & B\\
    0 & -A
\end{bmatrix}.$ The integrability condition of $\phi$ is that $A\wedge B = 0$. Irrespective of $B=0$ or $B\neq 0$ it follows that $\det(\phi) = -\mbox{Sym}^2(A)$.\\

Considering all possible cases of the determinants, it follows that \begin{align*}
&\mbox{Im}\left(\det|_{\mathcal{M}_{\mathbb{P}^2}(V_1^\rho, T_{\mathbb{P}^2})}\right)\\
&= \left\{q\otimes\mbox{Sym}^2(C): q\in H^0(\mathcal{O}(2)),~ C\in H^0(T_{\mathbb{P}^2}(-1))\right\}\bigcup \left\{\mbox{Sym}^2(A): A\in H^0(T_{\mathbb{P}^2})\right\}\\
&= \left\{q\otimes\mbox{Sym}^2(C): q\in H^0(\mathcal{O}(2))^\times,~ C\in H^0(T_{\mathbb{P}^2}(-1))^\times\right\}\bigcup \left\{\mbox{Sym}^2(A): A\in H^0(T_{\mathbb{P}^2})\right\}.
\end{align*}

However, the intersection of these two collections of determinant sections in the last step is nonempty. There are sections $q\in H^0(\mathcal{O}(2)),~ C\in H^0(T_{\mathbb{P}^2}(-1))$ and $A\in H^0(T_{\mathbb{P}^2})$ such that $$q\otimes\mbox{Sym}^2(C) = \mbox{Sym}^2(A).$$ Such sections should be identified as the same. On the other hand sections of the form $q\otimes\mbox{Sym}^2(C)$ for $q\neq 0,~ C\neq 0$ are identified with points in $\frac{H^0(\mathcal{O}(2))^\times\times H^0(T_{\mathbb{P}^2}(-1))^\times}{(q, C)\sim (\alpha^2\cdot q, \alpha^{-1}\cdot C)}$ by $q\otimes\mbox{Sym}^2(C)\mapsto [(q, C)]$ and the sections of the form $\mbox{Sym}^2(A)$ are identified with points in $H^0(T_{\mathbb{P}^2})$. Under the identification $q\otimes\mbox{Sym}^2(C) = \mbox{Sym}^2(A)$ we arrive at a bijective correspondence 

\begin{align*}
&\mbox{Im}\left(\det|_{\mathcal{M}_{\mathbb{P}^2}(V_1^\rho, T_{\mathbb{P}^2})}\right)\to\\
&\left(\frac{H^0(\mathcal{O}(2))^\times\times H^0(T_{\mathbb{P}^2}(-1))^\times}{(q, C)\sim (\alpha^2\cdot q, \alpha^{-1}\cdot C)}\sqcup\frac{H^0(T_{\mathbb{P}^2})}{\{\pm 1\}}\right)\bigg/[(q, C)]\sim [A]\iff q\otimes\mbox{Sym}^2(C) = \mbox{Sym}^2(A).
\end{align*}

via the correspondence $q\otimes\mbox{Sym}^2(C)\mapsto [(q, C)];~~\mbox{Sym}^2(A)\mapsto [A]$.
\begin{remark}
Now we have
\begin{align*}
&\dim \mbox{Im}\left(\det|_{\mathcal{M}_{\mathbb{P}^2}(V_1^\rho, T_{\mathbb{P}^2})}\right) = \max \left\{\dim \frac{H^0(\mathcal{O}(2))^\times\times H^0(T_{\mathbb{P}^2}(-1))^\times}{(q, C)\sim (\alpha^2\cdot q, \alpha^{-1}\cdot C)};~\dim \frac{H^0(T_{\mathbb{P}^2})}{\{\pm 1\}}\right\}\\
& < \max\{9, 8\} = 9 < 27 = \dim H^0(\mbox{Sym}^2(T_{\mathbb{P}^2})).
\end{align*} Thus $\det|_{\mathcal{M}_{\mathbb{P}^2}(V_1^\rho, T_{\mathbb{P}^2})}$ is not surjective onto $H^0(\mbox{Sym}^2(T_{\mathbb{P}^2}))$.
\end{remark}



\end{proof}

\end{document}
